\newpage
\thispagestyle{empty}
\begin{spacing}{1.2}
\chapter{Problem Statement and Objectives}

This chapter defines the core problem that GramaGIS addresses and outlines the specific objectives guiding its design and implementation. Together, these elements establish the technical and functional scope of the project.

\section{Problem Statement}
To develop a web-based Smart Panchayath GIS portal that integrates fragmented rural administrative data into a spatial platform for transparent governance and efficient infrastructure management.

\section{Objectives}
The primary objective of GramaGIS is to design and implement a Smart Panchayat GIS Portal that centralizes village infrastructure data into an interactive and intelligent digital environment. The framework aims to simplify geospatial queries through natural language, making spatial data accessible to everyone.

The specific objectives are:
\begin{itemize}
\item \textbf{Spatial Digitization}: Digitize and organize village assets such as ward boundaries, schools, hospitals, and road networks within a standardized geospatial framework.
\item \textbf{Centralized Data Repository}: Develop a PostGIS-based spatial database to support efficient storage, relational management, and spatial querying of rural datasets.
\item \textbf{NLP Integration}: Integrate a natural language interface using the Gemini API that enables users to query the panchayat database using plain English (e.g., "Show all hospitals in Ward 5") and obtain visual or tabular results.
\item \textbf{Interactive Visualization}: Create a dynamic web-based map interface using Leaflet.js that allows users to visualize, explore, and filter village data layers in real-time.

\end{spacing}